\documentclass[12pt]{exam}

\usepackage{amsmath, amssymb}
\usepackage{graphicx}
\usepackage{enumitem}

\begin{document}

{\small Seed: 65877 \hfill Version 4}

\vspace{1em}

\begin{questions}



% ==================================================
% GROUPED QUESTIONS (parts)
% ==================================================


\question 

\begin{parts}


\part Solve the system: $y = 2x + 4, \\ 3x + y = 18$




\vspace{2\baselineskip}








\part Solve for $x$ and $y$: $2x + 3y = 20, \\ 4x - y = 9$




\vspace{2\baselineskip}







\end{parts}

% ==================================================
% REGULAR QUESTIONS (no grouping)
% ==================================================


\vspace{1em}



% ==================================================
% GROUPED QUESTIONS (parts)
% ==================================================


\question The determinant of a diagonal matrix equals the product of its diagonal entries. from Q5




\begin{itemize}[label={}]
  \item (A) True
  \item (B) False
\end{itemize}









\vspace{1em}



% ==================================================
% GROUPED QUESTIONS (parts)
% ==================================================


\question If two matrices $A$ and $B$ commute, then $AB = BA$.




\begin{itemize}[label={}]
  \item (A) True
  \item (B) False
\end{itemize}









\vspace{1em}



% ==================================================
% GROUPED QUESTIONS (parts)
% ==================================================


\question Answer True or False for each of the following statements:

\begin{parts}


\part The transpose of a symmetric matrix is the negative of itself.




\begin{itemize}[label={}]
  \item (A) True
  \item (B) False
\end{itemize}









\part The determinant of a triangular matrix is equal to the product of its diagonal entries.




\begin{itemize}[label={}]
  \item (A) True
  \item (B) False
\end{itemize}








\end{parts}

% ==================================================
% REGULAR QUESTIONS (no grouping)
% ==================================================


\vspace{1em}



% ==================================================
% GROUPED QUESTIONS (parts)
% ==================================================


\question Which of the following represents an antiderivative of \( e^{2x} \)?




\begin{choices}

  
    \choice $x e^{2x}$
  

  
    \CorrectChoice $\tfrac{1}{2} e^{2x}$
  

  
    \choice $2e^{2x}$
  

\end{choices}









\vspace{1em}



% ==================================================
% GROUPED QUESTIONS (parts)
% ==================================================


\question Donner tous les choix qui sont 1.




\begin{choices}

  
    \CorrectChoice [$1$,un,$e^0$]
  

  
    \choice $-1$
  

  
    \choice $\tfrac{\pi}{2}$
  

\end{choices}









\vspace{1em}



% ==================================================
% GROUPED QUESTIONS (parts)
% ==================================================


\question If \( \int_0^a f(x)\,dx = 5 \), what is \( \int_a^0 f(x)\,dx \)?




\begin{choices}

  
    \CorrectChoice $-5$
  

  
    \choice $10$
  

  
    \choice $0$
  

\end{choices}









\vspace{1em}



\end{questions}

\end{document}